\chapter{Experiment and Results}

\section{Introduction}

\hs This chapter presents the experimental results obtained from training and evaluating the crack segmentation models described in Chapter 3. The experiments were designed to address three primary research questions:
\begin{enumerate}
    \item How do different architectural paradigms (semantic segmentation, instance segmentation, and transformer-based models) compare in terms of crack detection accuracy and computational efficiency?
    \item Does the tile-based inference pipeline with blended overlap merging effectively preserve crack continuity when processing large-scale images?
    \item How does fine-tuning on domain-specific data affect model performance?
\end{enumerate}
The chapter begins by detailing the experimental environment and hardware specifications, followed by a summary of the dataset splits used for training, validation, and testing. Results are then presented for each model architecture, including performance metrics on public datasets, inference pipeline evaluations, and ablation studies investigating the contribution of individual components.

\newpage

\section{Experimental Environment and Hardware Specifications}

\hs All experiments were conducted using a consistent hardware and software environment to ensure fair comparison between models. The technical specifications are summarized below:

\begin{table}[htbp]
    \centering
    \caption{Experimental hardware and software specifications}
    \label{tab:hardware_specs}
    \begin{tabular}{l p{0.6\textwidth}}
        \toprule
        \textbf{Component}       & \textbf{Specification}                                            \\
        \midrule
        CPU                      & Intel i5-12600K (12 cores, 16 threads)                            \\
        RAM                      & 64 GB DDR4                                                        \\
        GPU                      & NVIDIA RTX A6000 (48 GB VRAM)                                     \\
        Operating System         & Ubuntu 22.04 LTS                                                  \\
        Python Version           & 3.9.18                                                            \\
        PyTorch Version          & 2.0.1+cu118                                                       \\
        CUDA Version             & 11.8                                                              \\
        cuDNN Version            & 8.9.0                                                             \\
        Deep Learning Frameworks & Segmentation Models PyTorch (SMP) v0.3.3, Ultralytics YOLO v8.2.0 \\
        Image Processing         & OpenCV 4.8.0, Albumentations 1.3.1                                \\
        \bottomrule
    \end{tabular}
\end{table}

\hs \hl{The NVIDIA RTX A6000 GPU with 48 GB of VRAM was selected to accommodate large batch sizes during training and to enable inference on high-resolution images through the tile-based pipeline.}

\section{Results from Baseline Architectures}
\subsection{Performance on Public Datasets}
\hs Each model was trained using the protocol described in \hl{Section 3.4} and evaluated on the held-out test split (882 images). Performance metrics are reported at the optimal threshold determined via validation set analysis.

\begin{table}[htbp]
    \centering
    \caption{Segmentation performance on test split}
    \label{tab:segmentation_performance}
    \begin{tabular}{p{4.5cm}ccccc}
        \toprule
        \textbf{Model}          & \textbf{IoU} & \textbf{P} & \textbf{R} & \textbf{F1} & \textbf{Time (ms)} \\
        \midrule
        U-Net (EfficientNet-B0) & 0.5304 ± 0.1711         & 0.6532 ± 0.1641       & 0.7316 ± 0.1869            & 0.6759 ± 0.1566        & 84.45 ± 12.18                 \\
        SegFormer (MiT-B0)      & 0.5121 ± 0.1816         & 0.6467 ± 0.1590       & 0.7084 ± 0.2130           & 0.6570 ± 0.1718        & 75.96 ± 16.19                 \\
        YOLO26s-seg             & 70.8         & 79.5       & 87.1            & 83.1        & 22                 \\
        \bottomrule
    \end{tabular}
\end{table}
\newpage

\subsection{Analysis of the Blended Overlap Approach}

\section{Results and Analysis of Fine-Tuning on Self-Collected UAV Data}
\subsection{Model Adaptation and Performance}
\subsection{Generalization to Unseen Large-Scale Data}

\section{Ablation Studies}

\section{Summary}